\documentclass[../../main.tex]{subfiles}
\begin{document}

\begin{flushright}
\footnotesize\textit{【自動文字起こし・要確認】}
\end{flushright}

\begin{proof}[解]
出発点の座標を $0$ とすると時刻 $t$ での A, B の座標は各々
\[
  A \dots a t
\]
\[
  B \dots v\left(t - \dfrac{l}{a}\right) \quad \left(t \ge \dfrac{l}{a}, v > a\right)
\]
とおける．B が A に追いつく時刻 $t_0$ は
\[
  t_0 = \dfrac{l/a}{1 - \dfrac{a}{v}} = \dfrac{l v}{a(v-a)}
\]
である．したがって，B のひろう $f(v)$ は
\begin{align}
  f(v) &= v^2 \left(t_0 - \dfrac{l}{a}\right) \\
  &= \dfrac{v^2}{v-a} l = l \left(v + a + \dfrac{a^2}{v-a}\right)
\end{align}
\[
  \dfrac{f'(v)}{l} = 1 - \dfrac{a^2}{(v-a)^2} = \dfrac{v(v-2a)}{(v-a)^2}
\]
以下表を考える．

\[
\begin{array}{c|c*3{c}}
v & a & \dots & 2a & \dots \\ \hline
f' & & - & 0 & + \\ \hline
f & & \searrow & & \nearrow
\end{array}
\]

よって $v = 2a$ の時，疲労が最小．
\end{proof}

\newpage

\end{document}
